\section{Experiment}
\label{sec:experiment}

% Budget: about half a page. The table takes ~8 lines, leaving 2-3 lines per part.
% This section turns the objects of Section 2 into procedures: p_test (unit tests),
% the swarm law (swarm runs), p_swarm(i,t) (replay), and L_model (statistics).

% Run matrix, from the hand-off. "Runs" is per cell, i.e. per (step, n, graph, k):
% this is what the hand-off's Modal cost estimate assumes (e.g. the core sweep:
% 600k calls x 650 tokens at 5,000 tokens/s = 21.7 GPU-hours). Say so in the caption.
% Caption, still to write: T = 20 rounds throughout; what one run is; that the
% conditional extras run only if Saturday is clean, in the order numbered.
% To save space, the three conditional rows can collapse into one line of the caption.
\begin{table}[ht]
  \centering\footnotesize
  \setlength{\tabcolsep}{4pt}
  \begin{tabular}{@{}p{0.15\textwidth}llllll@{}}
    \toprule
    Arm & Steps & $n$ & Graphs & $k$ & Runs & Status \\
    \midrule
    Calibration (last $k$ rounds) & Counts, cheap talk & 4          & Chain, star, fully connected & 1, 2 & 100          & Committed \\
    Core (full history)           & Counts, cheap talk & 10, 20, 40 & Chain, star, fully connected & 1    & 100, 100, 50 & Committed \\
    Scale                         & Counts, cheap talk & 100        & Cycle, random 4-regular      & 1    & 50           & Committed \\
    \midrule
    Free text                     & Free text          & 10, 20, 40 & Cycle, random 4-regular      & 1    & 50           & Conditional 1 \\
    Second 7B family              & Cheap talk         & 10, 20, 40 & Chain, star, fully connected & 1    & 100, 100, 50 & Conditional 2 \\
    $\sim$70B model               & Counts, cheap talk & 10         & Chain, star, fully connected & 1    & 50           & Conditional 3 \\
    \bottomrule
  \end{tabular}
  \caption{Run matrix: $T = 20$ rounds; one round follows Table~\ref{tab:ladder}; estimated time to run is 45 GPU hours; conditional extras run if Saturday is clean, in given order}
  \label{tab:runs}
\end{table}

% UNIT-TEST PROTOCOL (how p_test is built; no swarm data is used here: say so).
%  - Test per in-degree class, not per agent: K_i depends on i only through d_i, so
%    one set of tests per distinct d_i (chain: 1, 2; star: 1, n-1; fully connected:
%    n-1; cycle: 2; random 4-regular: 4). This is why unit tests are affordable.
%  - For each coarse state s, ~20 synthetic histories h with pi_k(q(h)) = s.
%    Lengths uniform over the lengths that can produce s: the empty state () only
%    from length 0, every other state needs length >= 1. Make this agree with the
%    "uniform on {0, ..., T-1}" in Definition 2.
%  - Earlier rounds random: actions, counts, and signals on the cheap-talk step;
%    each round's payoff computed from them, so unit-test prompts have exactly the
%    same format as swarm prompts.
%  - Exact probabilities: action token before signal in the output format, so
%    pr_*K_i is the first token's probability from one forward pass. State what
%    happens to mass on invalid tokens (e.g. renormalise over H and L).
%  - kappa_i(. | s) = average over the histories for s; eps_i(p_test(i,t)) = mean
%    TV distance over histories of length t-1.
We test per in-degree class instead of per agent.
As $K_i$ depends on $i$ only through $d_i$, we have one set of tests per distinct $d_i$.
This also results in unit tests being affordable.
For each coarse state $s$, we have about 20 synthetic histories $h$ with $\pi_k(q(h)) = s$.
History lengths are uniform over the lengths that can produce $s$.
The empty state $()$ takes from length $0$, and every other state needs length at least one.
Each round's payoff is computed from that round's own action and the neighbour action counts, so the unit-test prompts have the same format as the swarm prompts.
The action token is before the signal in the output format, so $\mathrm{pr}_* K_i$ is the first token's probability from one forward pass.
For invalid tokens, we renormalize over $\mathrm{H}$ and $\mathrm{L}$.
The average over the histories for $s$ is given by $\kappa_i(\cdot \mid s)$, and the mean total-variation distance over histories of length $t-1$ is given by $\varepsilon_i(\ptest(i,t))$.


% SWARM RUNS (the swarm law of Definition 1).
%  - Temperature 1; actions and signals sampled each round; all agents' calls go to
%    the server as one batch per round (the synchronous update).
%  - Record the full outputs o^t, not just actions, so every history can be rebuilt
%    through rho_i; replay depends on this.
%  - Prompt conditions from Section 2 (point back, don't repeat): payoff shown, no
%    running total, T undisclosed, round 1 is the prompt alone, calibration prompts
%    cut to the last k rounds.
We choose a temperature of one, our actions and signals are sampled each round, and all agents' calls go to the server as a single batch per round.
We record the full outputs $o^t$, so every history can be obtained via $\rho_i$; this is required for replay.
Prompt conditions may be found in Section~\ref{sec:set-up}.

% REPLAY (how p_swarm(i,t) is estimated).
%  - ~5,000 recorded histories per (size, topology, step), each fed to a single
%    agent; exact probabilities; TV distance from kappa_i.
%  - Stratify the sample by round, so eps_i(p_swarm(i,t)) exists for every t:
%    Theorem 1 sums over t, and H2 compares round by round.
%  - Swarm data only scores the model and never fits it: say so (a reviewer will
%    check this; it is in the hand-off's risk table).
Our sample size is approximately 5000 recorded histories per size, topology, and step; they are each fed into a single agent.
We have exact probabilities, and the total-variation distance is obtained from $\kappa_i$.
We stratify the sample by round, so that $\varepsilon_i(\pswarm(i,t))$ exists for every $t$.
Theorem~\ref{thm:TV_bound} sums over $t$, and H2 compares each round.
The swarm data only scores the model; it does not fit the model.

% STATISTICS AND DELTA (L_model against the swarm).
%  - Statistics: fraction bidding high per round; time to consensus; final split.
%    Define consensus precisely (e.g. all agents take the same action for m
%    consecutive rounds), or time to consensus cannot be reproduced.
%  - Model side: Monte Carlo of the composed chain, with many more runs than the
%    swarm (it is cheap); exact for n = 4.
%  - Delta: TV distance between swarm and model distributions of each statistic;
%    bootstrap 95% intervals across runs. Bin statistics with many values: TV
%    between empirical distributions is biased upwards with many bins.
%  - Plot Delta against T with two curves: Theorem 1's bound with eps_i(p_swarm),
%    and the same sum with eps_i(p_test). Label the second a prediction, not a
%    bound: the theorem needs p_swarm.
%  - Short horizon: T = 5, from the first five rounds of each run (one clause).
Note that total variation, which is bounded by Theorem~\ref{thm:TV_bound}, between the two laws of a statistic, is a lower bound on the true trajectory distance (which is at most the bound), so we may use statistics to test our theorem.
We define $X_t = \tfrac{1}{n} \sum_i \mathbf{1}[a^t_i = \mathrm{H}]$, which gives one $\Delta_t$ per round; this is what the $\Delta$ versus $T$ plot needs.
The first round from which all agents take the same action for 3 consecutive rounds is denoted by $\tau$.
If we have no consensus, then we write $\tau = \infty$.
The final split is the average of $X_t$ over the last 5 rounds.

We have a Monte Carlo simulation of the composed chain, where each step only looks up $\kappa_i$, so $10^4$--$10^5$ runs are cheap, and the model's sampling noise becomes negligible.
For $n=4$, we may compute it exactly.
We estimate total variation from 100 samples, with the estimated total variation being biased upwards.
With $n=40$, $X_t$ takes 41 values, and two samples drawn from the same law can still give a total variation of $0.2$--$0.3$.
We thus bin $X_t$ into five bins of width $0.2$.
We measure the noise floor by taking two samples of the swarm's size from the model, compute the total variation between them, and plot it as a band.
Also, for a statistic with values in $[0, 1]$, $\lvert \mathbb{E}_{\text{swarm}} X_t - \mathbb{E}_{\text{model}} X_t \rvert \leq \Delta_t$, where we are considering unbinned $\Delta_t$, so we report the difference of means next to the binned total variation.
We bootstrap 95\% intervals across runs, resampling whole runs with replacement and recomputing $\Delta$.

Our plot of $\Delta$ against $T$ has two curves.
One is Theorem~\ref{thm:TV_bound}'s bound $B(t) = \sum_{s \leq t} \sum_i \varepsilon_i(\pswarm(i, s))$, and the other is the same sum with $\varepsilon_i(\ptest(i, s))$, which is a test-only prediction.

% BASELINES.
%  - Order-0 audit: kappa_i^0 = agent i's action distribution under p_test,
%    ignoring its coarse state; the composed chain then has independent agents.
%    Expected to fail, as in Ashery et al. (add to refs.bib).
%  - Order-1 audit: the method.
%  - Markov chain fitted to swarm data: the upper benchmark, as in Hirota et al.
%    (add to refs.bib). Say it is not an audit, since it uses swarm data. Fit on
%    half the runs and evaluate on the other half, so it cannot win by memorising.
% \placeholder{Baselines}

% The Friday pilot belongs with Section 6's milestones, not here.
